\documentclass[aspectratio=169,10pt]{beamer}

%% Use the official HUST theme
%% Options:
%%   color       = blue | red
%%   font        = lato | arial | times
%%   titlefont   = same | avantgarde
%%   footerlogo  = full | mark | wordmark | hust | none
%%   bar         = thin | thick
%%   titlepage   = cover | dark
%%   sectionpage = none | panel
%%   numbering   = none | plain | fraction
%%   footer      = <custom footer text>
%%   footerleft  = <custom left text>
\usetheme[
  color=blue,
  font=cm,
  titlefont=same,
  footerlogo=mark,
  bar=thin,
  titlepage=cover,
  sectionpage=panel,
  numbering=fraction,
  progressbar=on,
  logo2image=fami-shield.png,
  logo2vn={KHOA TOÁN - TIN},
  logo2en={Faculty of Mathematics and Informatics},
  logo2color=famiblue
]{HUST}

\usepackage{amsmath,amssymb}
\usepackage{booktabs}
\usepackage[backend=biber,style=numeric,sorting=none,maxcitenames=2,mincitenames=1]{biblatex}
\addbibresource{refs.bib}

\title{Xây dựng mẫu trình chiếu Beamer chuẩn nhận diện HUST}
\subtitle{Hệ thống trình chiếu LaTeX hiện đại, linh hoạt tỉ lệ và kích thước phông chữ}
\author{Nguyễn Văn A}
\institute{Trường Công nghệ Thông tin và Truyền thông\\Đại học Bách khoa Hà Nội}
\date{\today}

\begin{document}

%% ------------------------------------------------------------- 0a. Welcome Screen 1: Normal bg-dark
\begin{hustwelcomeframe}
\end{hustwelcomeframe}

%% ------------------------------------------------------------- 0b. Welcome Screen 2: Modified bg-background showcase
\begin{hustwelcomeframe}[bg=background, footer={DEMO SHOWCASE}]
\end{hustwelcomeframe}

%% ------------------------------------------------------------- 1. Title Slide (Image 1)
\hustsetup{slogantext={DEMO PRESENTATION}}
\begin{frame}[plain,noframenumbering]
  \titlepage
\end{frame}

%% ------------------------------------------------------------- 2. Outline
\begin{frame}{\husttoctitle}{\husttocsubtitle}
  \tableofcontents[split=3,hideallsubsections]
\end{frame}

%% ------------------------------------------------------------- Section 1
\section{Giới thiệu \& quy chuẩn}
\subsection{Hệ thống nhận diện thương hiệu}
\begin{frame}{Tạo cụm logo động theo sổ tay thương hiệu}{Quy chuẩn Hình 7 (Việt - Anh) và Hình 9 (Đơn vị thành viên)}
  \scriptsize
  \begin{columns}[T]
    \begin{column}{0.48\textwidth}
      \textbf{1. Cụm chính thức Việt - Anh (Hình 7):}
      \vskip0.3em
      \hustlogo[height=0.85cm]
      \vskip0.7em
      \textbf{2. Bản 2 dòng Việt - Anh (Hình 7b):}
      \vskip0.3em
      \hustlogo[height=0.85cm, layout=2, vn={ĐẠI HỌC\newline BÁCH KHOA HÀ NỘI}, en={HANOI UNIVERSITY\newline OF SCIENCE AND TECHNOLOGY}]
      \vskip0.7em
      \textbf{3. Đơn vị thành viên 1 dòng (Hình 9):}
      \vskip0.3em
      \hustlogo[height=0.85cm, unit={TRƯỜNG CƠ KHÍ}]
    \end{column}

    \begin{column}{0.48\textwidth}
      \textbf{4a. Đơn vị 2 dòng -- ĐHBKHN 2 dòng (Đối xứng):}
      \vskip0.3em
      \hustlogo[height=0.85cm, layout=4a, unit={VIỆN CÔNG NGHỆ ĐIỀU KHIỂN\newline VÀ TỰ ĐỘNG HÓA}]
      \vskip0.7em
      \textbf{4b. Đơn vị 2 dòng -- ĐHBKHN 1 dòng (Căn giữa):}
      \vskip0.3em
      \hustlogo[height=0.85cm, layout=4b, unit={VIỆN CÔNG NGHỆ ĐIỀU KHIỂN\newline VÀ TỰ ĐỘNG HÓA}]
      \vskip0.7em
      \textbf{5. Cụm logo Khoa Toán - Tin (FaMI, Hình 3):}
      \vskip0.3em
      \hustlogo[image=fami-shield.png, height=0.85cm, vn={KHOA TOÁN - TIN}, en={Faculty of Mathematics and Informatics}, color=famiblue, equalwidth=false]
    \end{column}
  \end{columns}
\end{frame}

\subsection{Quy định màu sắc và phông chữ}

\begin{frame}{Hệ thống nhận diện thương hiệu ĐHBK Hà Nội}{Tuân thủ Sổ tay nhận diện thương hiệu và mẫu PowerPoint 2026}
  \begin{itemize}
    \item \textbf{Màu sắc quy chuẩn}:
      \begin{itemize}
        \item Hệ màu thứ cấp: Xanh lam đậm (\texttt{\#003366}), Đỏ (\texttt{\#C02034}).
        \item Hệ màu chính: Đỏ cờ (\texttt{\#CE1628}), Vàng sao (\texttt{\#F3C108}), Đen (\texttt{\#231F20}).
      \end{itemize}
    \item \textbf{Kiểu chữ quy định}:
      \begin{itemize}
        \item Phông chữ chính (tiêu đề): \textbf{Avant Garde} (TeX Gyre Adventor).
        \item Phông chữ phụ (nội dung): \textbf{Lato}.
        \item Phông chữ văn bản: \textbf{Times New Roman}.
      \end{itemize}
    \item \textbf{Tỉ lệ khung hình linh hoạt}: Hỗ trợ tự nhiên 16:9, 4:3, 16:10, 5:4 qua tùy chọn lớp \texttt{aspectratio}.
  \end{itemize}
\end{frame}

%% ------------------------------------------------------------- Section 2
\section{Các khối nội dung \& toán học}
\subsection{Các loại khối hiển thị}

\begin{frame}{Các loại khối hiển thị (Beamer blocks)}{Đồng bộ màu sắc nhận diện thương hiệu}
  \begin{block}{Khối tiêu chuẩn (hustmain)}
    Khối định nghĩa hoặc thông tin quan trọng được làm nổi bật với màu xanh nhận diện của trường.
  \end{block}

  \begin{alertblock}{Khối cảnh báo (hustalert)}
    Khối cảnh báo hoặc kết luận then chốt sử dụng màu đỏ thương hiệu HUST.
  \end{alertblock}

  \begin{exampleblock}{Khối ví dụ (hustorange)}
    Ví dụ minh họa, mã nguồn hoặc kết quả thực nghiệm với màu cam nhận diện HUST.
  \end{exampleblock}
\end{frame}

\subsection{Công thức toán học \& bảng số liệu}
\begin{frame}{Công thức toán học \& bảng dữ liệu}{Hiển thị sắc nét ở mọi độ phân giải}
  \begin{columns}[T]
    \begin{column}{0.48\textwidth}
      \textbf{Mô hình tối ưu hóa}:
      \begin{equation}
        \min_{\mathbf{w}} \frac{1}{2}\|\mathbf{w}\|^2 + C \sum_{i=1}^N \xi_i
      \end{equation}
      thỏa mãn điều kiện:
      \begin{equation}
        y_i (\mathbf{w}^T \mathbf{x}_i + b) \ge 1 - \xi_i, \quad \xi_i \ge 0
      \end{equation}
    \end{column}

    \begin{column}{0.48\textwidth}
      \textbf{Kết quả thực nghiệm}:
      \begin{table}
        \centering
        \small
        \begin{tabular}{lrr}
          \toprule
          \textbf{Mô hình} & \textbf{Độ chính xác} & \textbf{F1-Score} \\
          \midrule
          Baseline       & 87.5\% & 0.86 \\
          CNN            & 93.2\% & 0.92 \\
          Transformer    & \textbf{96.8\%} & \textbf{0.96} \\
          \bottomrule
        \end{tabular}
      \end{table}
    \end{column}
  \end{columns}
\end{frame}

%% ------------------------------------------------------------- Section 3: Biber / BibLaTeX Citations
\section{Trích dẫn \& tài liệu tham khảo}
\subsection{Phương pháp tối ưu hóa đa mục tiêu}

\begin{frame}{Nghiên cứu tối ưu hóa đa mục tiêu}{Trích dẫn 1 dòng: thanh footer giữ nguyên vị trí phân tách tinh tế}
  Nghiên cứu về phương pháp tối ưu hóa đa mục tiêu được đề xuất bởi Bello-Cruz và cộng sự\cite{bellocruz2025} để giải Bài toán (4).

  \begin{itemize}
    \item Nếu $C$ là một tập compact, theo Định lý 6\cite{bellocruz2025}, dãy $\{x_k\}$ sinh ra bởi thuật toán MPG-E hội tụ về một nghiệm Pareto yếu của bài toán (4).
    \item Cụm trích dẫn tự động căn lề tránh logo chân trang (\texttt{footerlogo=mark}), thanh viền footer tự động điều chỉnh linh hoạt.
    \item Trích dẫn lặp lại cùng tài liệu trên cùng slide sẽ tái sử dụng chỉ số mà không bị trùng lặp ở chân trang.
  \end{itemize}
\end{frame}

\subsection{Thuật toán quán tính trong không gian Hilbert}
\begin{frame}{Nhiều trích dẫn trên cùng một slide}{Thanh footer tự động đẩy lên 3 dòng để chứa trích dẫn}
  Các thuật toán giải bài toán tối ưu và ứng dụng trong khoa học kỹ thuật:
  \begin{itemize}
    \item Phương pháp gradient gần kề với tìm kiếm đường thẳng rõ ràng\cite{bellocruz2025}.
    \item Thuật toán quán tính cho bài toán tối ưu lồi trong không gian Hilbert\cite{alvarez2024}.
    \item Xuất bản khoa học và chuẩn nhận diện thương hiệu đại học\cite{hust2026}.
  \end{itemize}
\end{frame}

%% ------------------------------------------------------------- Section 4 (Tắt panel slide để minh họa tùy chọn)
\hustsection[panel=false]{Tùy biến \& linh hoạt}

\begin{frame}{Các tùy chọn của gói theme}{Tùy biến dễ dàng qua \texttt{\textbackslash usetheme} hoặc \texttt{\textbackslash hustsetup}}
  \begin{description}
    \item[\texttt{color}] \texttt{blue} (mặc định) hoặc \texttt{red}.
    \item[\texttt{font}] \texttt{lato} (mặc định), \texttt{arial}, \texttt{times}, hoặc \texttt{cm}.
    \item[\texttt{titlefont}] \texttt{same} hoặc \texttt{avantgarde}.
    \item[\texttt{slogan}] \texttt{on} (mặc định, hiển thị ONE LOVE. ONE FUTURE.) hoặc \texttt{off}.
    \item[\texttt{wordmark}] \texttt{on} (mặc định) hoặc \texttt{off} (bật/tắt cụm chữ HUST trung tâm).
    \item[\texttt{panelsocial}] \texttt{on} (mặc định) hoặc \texttt{off} (ẩn web/facebook trên panel).
    \item[\texttt{unit}] Tên trường/viện thành viên (ví dụ \texttt{unit=\{TRƯỜNG CƠ KHÍ\}}).
    \item[\texttt{footerlogo}] \texttt{full}, \texttt{mark}, \texttt{wordmark}, \texttt{hust}, hoặc \texttt{none}.
    \item[\texttt{numbering}] \texttt{plain}, \texttt{fraction}, hoặc \texttt{none}.
  \end{description}
\end{frame}

%% ------------------------------------------------------------- 5. References Frames (Manual Split with Identical Title)
\begin{hustreferencesframe}[from=1, to=2]
\end{hustreferencesframe}

\begin{hustreferencesframe}[from=3, to=3]
\end{hustreferencesframe}

%% ------------------------------------------------------------- Ending Dark Slide
\begin{hustdarkframe}
  \centering
  \vspace*{0.2\paperheight}
  {\Huge\bfseries\color{white} CẢM ƠN QUÝ THẦY CÔ\\VÀ CÁC BẠN ĐÃ LẮNG NGHE!\par}
  \vspace{1.5em}
  {\large\color{white!85} Q \& A\par}
\end{hustdarkframe}

\end{document}
